\section{Introduction}

\subsection{Context} % what the thesis is about
% This project explores the application of typed assembly languages in static
% worst-case execution time analysis.
% 
% In real-time systems the time it takes a program to run is per definition
% greatly important, since a real-time system is defined as being able to execute
% some code within a hard time limit.

ssd\_os is an operating system for solid state drives (ssd) equipped with RISC-V
processors, that provides a synchronous and deterministic programming model. A
program is a collection of pipelines; each pipeline is a linear sequence of
stages that execute to completion within a tick, i.e., a fixed period of time.
Ideally, the programming toolchain verifies that a given stage can execute
within a tick. This requires a static analysis of the RISC-V assembly program
to reason about the deterministic nature of a stage or its worst-case execution
time. In this project, we base our approach on typed assembler program. We rely
on the assembly type system to reason about determinism or running time.

\subsection{Problem} % the specific problem you focus on
In this project the problem was reduced to designing a subset of a typed
variant of risc-v to model a specific memory structure that used scratch and
shared memory regions, and to use this language and abstract machine model to
reason about the worst-case execution time of a program written in the typed
assembly language.

\subsection{Approach} % experimental, designing/implementing/evaluating a system
The approach was design heavy in the sense that the language and tool were new.
Initally, a lot of time was spent on considering in what way typed assembly (a
tool used for correctness) could be used as a tool for execution time analysis.

\subsection{Contribution} % an itemized list of what you have done
The contents of this project --- what has been done, can be summarized as the following
\begin{itemize}
    \item
        A brief presentation of the target system
    \item
        A high level overview of methods of worst-case execution time analysis
    \item
        An exploration of typed assembly
    \item
        How one may use typed assembly for worst-case execution time
    \item
        Designing a language based on the previous results and considerations
    \item
        and implementing a tool that analyses programs written in said language
        to reason about execution-time analysis.

\end{itemize}
